% ==============================================================================
% APÊNDICES DA DISSERTAÇÃO (PROPECAUT / IFES)
% ==============================================================================

\newpage
\phantomsection
\section*{APÊNDICES}
\addcontentsline{toc}{section}{APÊNDICES}

% ------------------------------------------------------------------------------
% APÊNDICE A
% ------------------------------------------------------------------------------
\phantomsection
\subsection*{APÊNDICE A -- Contrato Único de Parâmetros e Rastreabilidade de Fontes}
\addcontentsline{toc}{subsection}{APÊNDICE A -- Contrato Único de Parâmetros e Rastreabilidade de Fontes}

Este apêndice consolida o contrato único de parâmetros adotado nos módulos auditados (\texttt{contrato\_parametros.py}, conjunto \texttt{PROJETO\_1P0}), distinguindo explicitamente as grandezas verificadas em fontes primárias (normas, resoluções astronômicas e folhas de dados) das hipóteses de projeto declaradas pelo autor.

\begin{table}[htbp]
\centering
\small
\caption{Contrato único de parâmetros do nó orbital de $50\text{ kW}$ (\texttt{PROJETO\_1P0}).}
\label{tab:contrato_parametros}
\begin{tabularx}{\textwidth}{>{\raggedright\arraybackslash}p{4.8cm}cc>{\raggedright\arraybackslash}X}
\toprule
\textbf{Parâmetro} & \textbf{Valor} & \textbf{Unidade} & \textbf{Origem / Classificação} \\
\midrule
Semieixo maior $a_E$; $1/f$; $\mu_E$; $J_2$ & $6.378.137$; $298{,}257$; $3{,}986\times 10^{14}$; $1{,}0826\times 10^{-3}$ & SI & Fonte primária: NGA WGS84 \cite{nga2014wgs84} \\
Unidade astronômica; $R_\odot$; $S_0$ & $149.597.870.700$; $6{,}957\times 10^8$; $1.360{,}8$ & m; m; W/m$^2$ & Fontes primárias: IAU B2/B3 \cite{iau2012b2,iau2015b3}, Kopp \& Lean \cite{kopp2011tsi} \\
Albedo; emissão IR terrestre & $0{,}30 \pm 0{,}02$; $237 \pm 7$ & —; W/m$^2$ & Fonte primária: NASA SP-8067 \cite{nasa1971sp8067} \\
Órbita SSO ($h$, $i$, LTAN, $t_{ecl}$) & $700$; $98{,}188^\circ$; $18\text{h}00$; $1.105{,}39$ & km; $^\circ$; h; s & Calculado (\texttt{modelo\_orbital.py}); ref. \cite{nasa2018gdc} \\
Célula FV Spectrolab XTJ (BOL) & $29{,}5\%$; $26{,}62$; $2{,}34$; $17{,}0$ & \%; cm$^2$; V; mA/cm$^2$ & Fonte primária: Datasheet Spectrolab XTJ \cite{spectrolab_xtj_prime} \\
Arranjo FV EOL ($50\text{ kW}$ a $65^\circ\text{C}$) & $27.073$ cél.; $161{,}49$ / $179{,}43$ & —; m$^2$ & Calculado (\texttt{modelo\_fv.py}) \\
Eficiências ($\eta_{buck}$, $\eta_{dab}$, $\eta_{bd}$, $\eta_{pol}$) & $0{,}91$; $0{,}9655$; $0{,}97$; $1{,}00$ & — & Hipóteses de projeto (fronteira $48\text{ V}$) \\
Carga auxiliar ($P_{aux}$ / $Q_{aux}$) & $1.000{,}0$ / $903{,}2$ & W & Hipótese de projeto (fração óptica OISL) \\
Célula BESS Saft VES~16 & $3{,}6$; $4{,}5$; $0{,}115$; $140$ & V; Ah; kg; Wh/kg & Fonte primária: Datasheet Saft VES~16 \cite{saft_ves16} \\
Banco BESS ($19\text{P}108\text{S}$, $SoC$) & $2.052$ cél.; $236{,}0$; $[0{,}70,\; 0{,}95]$ & —; kg; — & Calculado (\texttt{modelo\_bess.py}); EOL $20\%$ (hip.) \\
Eficiências BESS ($\eta_{ch}$, $\eta_{dis}$) & $0{,}995$; $0{,}97$ & — & Hipóteses de projeto \\
Módulos IA (NVIDIA A100 SXM4) & $100 \times 400$; $624$ & W; TOPS & Fonte primária: Datasheet A100 \cite{nvidia_a100_datasheet} \\
Curva DVFS $V(f)$ e fuga ($\beta_T$, $\alpha_{leak}$) & $[300,\; 1.410]$; $0{,}025$; $0{,}15$ & MHz; K$^{-1}$; — & Hipóteses de projeto calibradas em $400\text{ W}$ \\
Planta térmica ($R_{int}$, $C_{mod}$, $C_{fluid}$) & $0{,}08$; $600$; $60.000$ & K/W; J/K; J/K & Hipóteses de projeto (\texttt{modelo\_termico.py}) \\
Revestimento Z-93 EOL (30 anos LEO) & $\alpha_s = 0{,}36$; $\varepsilon = 0{,}90$ & — & Resumo oficial NASA NTRS \cite{hasegawa1992z93} \\
Área de projeto do radiador ($A_{pl}$ / $A_{emit}$) & $49{,}90$ / $99{,}80$ & m$^2$ & Calculado para $38{,}18\text{ kW}$ ($T_j \le 80^\circ\text{C}$) \\
Sensores ($T_{raiz}$ RTD; $T_j$ diodo) & $\sigma=0{,}05$ ($\pm 0{,}25$); $\sigma=0{,}5$ ($\pm 0{,}5$/$1{,}5$) & K & IEC 60751 Classe A / Hipóteses declaradas \\
\bottomrule
\end{tabularx}
\end{table}

% ------------------------------------------------------------------------------
% APÊNDICE B
% ------------------------------------------------------------------------------
\newpage
\phantomsection
\subsection*{APÊNDICE B -- Catálogo de Módulos Auditados e Protocolo de Reprodutibilidade}
\addcontentsline{toc}{subsection}{APÊNDICE B -- Catálogo de Módulos Auditados e Protocolo de Reprodutibilidade}

A totalidade dos códigos-fonte, suítes de testes automatizados (\texttt{pytest}), checkpoints treinados, somas de verificação (\texttt{CHECKSUMS.txt}) e relatórios em JSON (UTF-8) encontra-se publicada sob licença livre no repositório oficial da dissertação (\url{https://github.com/dafnhe-thor-barros/mestrado}, pasta \texttt{simulacoes/v2\_auditado/}). A Tabela~\ref{tab:catalogo_modulos} sintetiza os módulos auditados e o placar de verificação local em Python 3.14 (NumPy, SciPy, Matplotlib e PyTorch 2.15 CPU).

\begin{table}[htbp]
\centering
\small
\caption{Módulos computacionais auditados e suítes de verificação em \texttt{simulacoes/v2\_auditado/}.}
\label{tab:catalogo_modulos}
\begin{tabularx}{\textwidth}{>{\ttfamily\raggedright\arraybackslash}p{4.6cm}c>{\raggedright\arraybackslash}X}
\toprule
\textbf{\textrm{Arquivo}} & \textbf{Testes} & \textbf{Responsabilidade e Escopo} \\
\midrule
modelo\_orbital.py & 24/24 & Órbita SSO 700 km WGS84, posição solar USNO, eclipse penumbral, fator de visão $F(90^\circ)$ \\
modelo\_fv.py & 11/11 & Dimensionamento EOL Spectrolab XTJ ($29{,}5\%$) sob AM0 para $50\text{ kW}$ \\
modelo\_barramento.py & 10/10 & Balanço HVDC $400\text{ V}$, Buck, DAB, bidirecional e fronteira de perdas \\
modelo\_bess.py & 10/10 & Dimensionamento Saft VES~16 ($108\text{S}$), dinâmica de $SoC \in [0{,}70,\; 0{,}95]$ \\
modelo\_ia.py & 11/11 & Modelo DVFS + fuga térmica para $100\times$ NVIDIA A100 SXM4 ($40\text{ kW}$ pico) \\
modelo\_termico.py & 12/12 & Planta de referência EDP 1D radiativa + EDO módulo (Crank-Nicolson/Rannacher + Newton) \\
modelo\_pinn.py & 14/14 & PINN v2 em PyTorch puro (\texttt{hard\_ic}, perdas em K-eq, Adam + L-BFGS) e estimador de $\Delta_{ponta}$ \\
auditoria\_marco4\_robusta.py & — & Auditoria independente em $500$ casos LHS, curva OOD, baselines RC2/RC3 e speedup \\
contrato\_parametros.py & — & Contrato único \texttt{PROJETO\_1P0}/\texttt{LEGADO\_2B\_0P90} e verificação de \texttt{SHA256} \\
marco5\_tems.py & 14/14 & Malha fechada TEMS/DVFS (S0, S1, R, RC3, PINN, Oráculo), barreira OOD e supervisor $SoC$ \\
\bottomrule
\end{tabularx}
\end{table}

\noindent\textbf{Total de verificações automatizadas:} $106$ testes unitários, de convergência, de conservação de energia e de integridade criptográfica (\texttt{SHA256}) executados com $100\%$ de aprovação local. O checkpoint congelado da PINN para o Marco~5 é \texttt{v2c.pt} (\texttt{SHA256: 05e32e16c6c0c7fb7f7486e418aef4bce7370c5f3d865726d00b4c2e7c29ff98}), com campanha multi-semente executada em GPU via \texttt{Treino\_PINN\_MultiSeed\_Colab.ipynb}.
